\documentclass[11pt]{article}

\usepackage[T1]{fontenc}
\usepackage[utf8]{inputenc}
\usepackage{lmodern}
\usepackage[margin=1.05in]{geometry}
\usepackage{microtype}
\usepackage{amsmath,amssymb,amsthm,mathtools}
\usepackage{bm}
\usepackage{booktabs}
\usepackage{longtable}
\usepackage{array}
\usepackage{enumitem}
\usepackage{xcolor}
\usepackage{hyperref}
\usepackage[nameinlink,noabbrev]{cleveref}
\usepackage{listings}
\usepackage{fancyhdr}

\hypersetup{
  colorlinks=true,
  linkcolor=blue!55!black,
  citecolor=green!40!black,
  urlcolor=blue!65!black,
  pdftitle={New structural progress on the Alaoglu--Erdos conjecture},
  pdfauthor={Research progress report}
}

\pagestyle{fancy}
\fancyhf{}
\lhead{Alaoglu--Erd\H{o}s progress report}
\rhead{24 August 2026}
\cfoot{\thepage}
\setlength{\headheight}{14pt}

\newtheorem{theorem}{Theorem}[section]
\newtheorem{proposition}[theorem]{Proposition}
\newtheorem{lemma}[theorem]{Lemma}
\newtheorem{corollary}[theorem]{Corollary}
\newtheorem{conjecture}[theorem]{Conjecture}
\newtheorem{problem}[theorem]{Problem}
\newtheorem{criterion}[theorem]{Criterion}
\newtheorem{observation}[theorem]{Observation}
\theoremstyle{definition}
\newtheorem{definition}[theorem]{Definition}
\newtheorem{experiment}[theorem]{Experiment}
\newtheorem{target}[theorem]{Target}
\theoremstyle{remark}
\newtheorem{remark}[theorem]{Remark}
\newtheorem{warning}[theorem]{Warning}

\newcommand{\AEconj}{\textsc{AE}}
\newcommand{\N}{\mathbb N}
\newcommand{\Z}{\mathbb Z}
\newcommand{\Q}{\mathbb Q}
\newcommand{\R}{\mathbb R}
\newcommand{\C}{\mathbb C}
\newcommand{\F}{\mathbb F}
\newcommand{\Gm}{\mathbb G_m}
\newcommand{\ord}{\operatorname{ord}}
\newcommand{\rank}{\operatorname{rank}}
\newcommand{\Res}{\operatorname{Res}}
\newcommand{\Norm}{\operatorname{N}}
\newcommand{\Tr}{\operatorname{Tr}}
\newcommand{\Gal}{\operatorname{Gal}}
\newcommand{\Span}{\operatorname{span}}
\newcommand{\im}{\operatorname{im}}
\newcommand{\supp}{\operatorname{supp}}
\newcommand{\content}{\operatorname{cont}}
\newcommand{\diag}{\operatorname{diag}}
\newcommand{\Li}{\operatorname{Li}}
\newcommand{\e}{\mathrm e}
\newcommand{\dd}{\,\mathrm d}
\newcommand{\ellr}{\ell}
\newcommand{\cL}{\mathcal L}
\newcommand{\cC}{\mathcal C}
\newcommand{\cR}{\mathcal R}
\newcommand{\cH}{\mathcal H}
\newcommand{\cQ}{\mathcal Q}
\newcommand{\cP}{\mathcal P}
\newcommand{\cK}{\mathcal K}
\newcommand{\cO}{\mathcal O}
\newcommand{\vct}[1]{\bm{#1}}
\newcommand{\abs}[1]{\left|#1\right|}
\newcommand{\norm}[1]{\left\lVert#1\right\rVert}
\newcommand{\angles}[1]{\left\langle#1\right\rangle}
\newcommand{\set}[1]{\left\{#1\right\}}
\newcommand{\floor}[1]{\left\lfloor#1\right\rfloor}
\newcommand{\ceil}[1]{\left\lceil#1\right\rceil}

\lstset{
  basicstyle=\ttfamily\small,
  columns=fullflexible,
  frame=single,
  breaklines=true,
  showstringspaces=false
}

\title{\textbf{New Structural Progress on the Alaoglu--Erd\H{o}s Conjecture}\\[0.35em]
\large A uniform staircase lattice, fixed-rank holonomy, and a diagonal Kummer supernorm}
\author{Research progress report continuing \texttt{alaoglu-erdos-unified-report-6}}
\date{24 August 2026}

\begin{document}
\maketitle

\begin{abstract}
The Alaoglu--Erd\H{o}s conjecture asserts that a real number $x$ satisfying
$2^x,3^x\in\Z$ must be an integer.  This report is a deliberately non-duplicative
continuation of the large unified report supplied with the problem.  It does
\emph{not} claim a proof of the conjecture.  It supplies two new exact pieces of
structure and uses them to sharpen the remaining research targets.

First, the earlier finite computations in the consecutive-log basis are upgraded
to a theorem valid for every prime smoothness cutoff $q$.  The map from the
$2\pi(q)$ prime-exponent coordinates of a candidate pair $(M,A)$ to the
$2q-3$ coefficients of its sparse weight-two period relation has kernel exactly
the one-dimensional control direction $(M,A)\mapsto(2^tM,3^tA)$, rank
$2\pi(q)-1$, and saturated image.  An explicit integral decoder recovers every
external valuation and the residual structural coordinate.  Equivalently, every
nonzero Smith invariant is $1$.  This turns the finite data observed at
$q=3,5,7,11$ into a uniform theorem and yields a fixed-rank, rank-six unipotent
connection whose length-two holonomy coefficient is precisely the
Alaoglu--Erd\H{o}s determinant.

Second, the elementary Kummer cocycle from the supplied report admits an exact
diagonal orbit product.  If $N\ge2$ and $P=2^{1/N}$, $Q=3^{1/N}$,
$U=M^{1/N}$, $V=A^{1/N}$, then with
$R=UQ-VP$ and $S=U+Q-V-P$ one has
\[
 \prod_{\zeta^N=1}
 \bigl((\zeta U-1)(\zeta Q-1)-(\zeta V-1)(\zeta P-1)\bigr)
 =R^N-S^N.
\]
At the distinguished real embedding this nonzero algebraic integer is not merely
$O(N^{-3})$: its absolute value is
$\exp(-N\log N+O_x(N))$, with an explicit asymptotic constant.  A complete
conjugate audit then explains why this striking gain is not yet a proof.  The
leading root-of-unity phase polynomial
\[
 F(p,q,u,v)=(u-1)(q-1)-(v-1)(p-1)
\]
has strictly positive logarithmic Mahler measure.  Thus the simple cocycle has
positive phase entropy, and multiplication by cyclotomic Kummer units cannot
remove it.  This gives a precise design condition for a successor cocycle: its
leading phase symbol must be generalized cyclotomic, while a later term must
still retain a nonzero external-prime descendant.

The report includes full proofs of the new lattice and asymptotic statements,
a reproducible numerical Mahler-measure computation, a Lean formalization
blueprint, and a prioritized set of sharply stated theorem targets.  The current
frontier is thereby narrowed from a broad collection of transcendence ideas to
two interfaces: restricted arithmetic holonomy on a saturated two-star residue
lattice, and construction of a control-sensitive Kummer cocycle with zero leading
phase entropy.
\end{abstract}

\tableofcontents
\newpage

\section{Problem, status, and scope}

\subsection{The conjecture and the counterexample package}

The conjecture under study is the classical question of Alaoglu and Erd\H{o}s~\cite{AE1944}.
\begin{conjecture}[Alaoglu--Erd\H{o}s]
If $x\in\R$ and $2^x,3^x\in\Z$, then $x\in\Z$.
\end{conjecture}

Assume for contradiction that a nonintegral counterexample exists and put
\begin{equation}\label{eq:MA}
 M=2^x\in\N,\qquad A=3^x\in\N.
\end{equation}
Then
\begin{equation}\label{eq:detzero}
 \log 2\,\log A-\log 3\,\log M=0.
\end{equation}
The rational case is elementary: if $x=r/s$ in lowest terms and $2^{r/s}$ is an
integer, unique factorization gives $s=1$.  The algebraic irrational case is
excluded by the Gelfond--Schneider theorem~\cite{Baker1975}.  Consequently any counterexample is
transcendental.  It is also positive, and current finite verification in the
supplied report places it beyond the range already checked there.

For all $n,k\in\N_0$, the same $x$ gives the semigroup of integer pairs
\begin{equation}\label{eq:semigrouppairs}
 2^{n+kx}=2^nM^k,\qquad 3^{n+kx}=3^nA^k.
\end{equation}
The supplied unified report~\cite{Unified6} develops this package from many viewpoints: valuation lattices,
compact irrational rotations, carry words, real and $p$-adic approximation,
Kummer extensions, arithmetic holonomy, matrix coefficients (including the recent advances of Dasgupta--Kakde~\cite{DasguptaKakdeII}), and small-type
entire interpolation.  The present document assumes that package and focuses on
claims not proved there.

\subsection{What is genuinely new here}

The supplied report already contains computational coefficient matrices for a
few small cutoffs and already identifies the finite weight-two period relation as
a plausible interface with the work of Calegari--Dimitrov--Tang.  It also already
constructs a small Kummer determinant and explains why a naively taken full field
norm loses the distinguished cancellation.  The following statements are the
new contributions of this continuation.

\begin{enumerate}[label=\textbf{N\arabic*.},leftmargin=3.6em]
\item \textbf{Uniform staircase theorem.}  The rank, kernel, saturation, Smith
normal form, and integral inverse of the consecutive-log coefficient map are
proved for every prime cutoff $q$, not merely checked by computer for selected
$q$.

\item \textbf{Prime-jump tomography.}  Every external prime valuation is an
explicit first difference of two residue coordinates.  The only invisible
integer direction is the control line.  This gives a complete local description
of the image lattice.

\item \textbf{Fixed-rank holonomy realization.}  The entire $q$-dependent
period relation is realized as one matrix coefficient of a rank-six unipotent
Fuchsian connection.  The number of punctures grows with $q$, but the rank does
not.

\item \textbf{Exact diagonal Kummer norm.}  The common root-of-unity orbit of
the elementary determinant is evaluated in closed form as $R^N-S^N$.

\item \textbf{Superexponential distinguished asymptotic.}  The exact leading
asymptotic of this algebraic integer is established.  Its logarithm is
$-N\log N+O_x(N)$.

\item \textbf{Phase-entropy obstruction.}  The leading torus symbol is proved
to have positive Mahler measure.  This turns the old qualitative warning about
``bad conjugates'' into a precise obstruction stable under cyclotomic-unit
renormalization.
\end{enumerate}

\subsection{Honest status}

None of the new theorems proves \AEconj.  The staircase theorem is a structural
advance: it replaces finite experimentation by a complete lattice theorem and
compresses the transcendence target to a fixed-rank holonomy problem.  The
Kummer theorem is an archimedean advance: it manufactures a much smaller
algebraic integer than the elementary cocycle.  The Mahler calculation then
shows exactly why the obvious norm completion fails.  A proof must add a new
mechanism at one of two places:
\begin{itemize}
\item a nonvanishing theorem for the restricted saturated coefficient lattice;
\item or a new Kummer cocycle whose leading phase symbol has zero Mahler measure
but whose first nontrivial descendant still detects an external valuation.
\end{itemize}

\section{The consecutive-log staircase}

\subsection{Telescoping coordinates}

Fix a prime $q\ge3$ and suppose that $M$ and $A$ are $q$-smooth.  Write
\begin{equation}\label{eq:valuations}
 M=\prod_{p\le q}p^{a_p},\qquad
 A=\prod_{p\le q}p^{b_p},
 \qquad a_p,b_p\in\N_0,
\end{equation}
where $p$ always denotes a prime.  For $1\le r\le q-1$ define
\begin{equation}\label{eq:Lr}
 L_r=\log\frac{r+1}{r}.
\end{equation}
The elementary telescoping identity
\begin{equation}\label{eq:logptelescope}
 \log p=\sum_{r=1}^{p-1}L_r
\end{equation}
turns prime valuations into tail counts.  Define
\begin{equation}\label{eq:tails}
 \alpha_r=\sum_{\substack{p\le q\\p>r}}a_p,
 \qquad
 \beta_r=\sum_{\substack{p\le q\\p>r}}b_p
 \qquad(1\le r\le q-1).
\end{equation}
Then
\begin{equation}\label{eq:logtail}
 \log M=\sum_{r=1}^{q-1}\alpha_rL_r,
 \qquad
 \log A=\sum_{r=1}^{q-1}\beta_rL_r.
\end{equation}
The sequences $(\alpha_r)$ and $(\beta_r)$ are nonnegative, nonincreasing,
integer-valued staircase functions.  Their jumps occur only at primes:
\begin{equation}\label{eq:jump}
 a_p=\alpha_{p-1}-\alpha_p,
 \qquad
 b_p=\beta_{p-1}-\beta_p,
\end{equation}
where $\alpha_q=\beta_q=0$ by convention.

The structural logarithms are
\begin{equation}\label{eq:structuralL}
 \log2=L_1,
 \qquad
 \log3=L_1+L_2.
\end{equation}
Substituting \eqref{eq:logtail} and \eqref{eq:structuralL} into
\eqref{eq:detzero} gives
\begin{equation}\label{eq:stairrelation}
 L_1\sum_{r=1}^{q-1}\beta_rL_r
 -(L_1+L_2)\sum_{r=1}^{q-1}\alpha_rL_r=0.
\end{equation}
This relation uses only the two stars through $L_1$ and $L_2$.

\subsection{The formal coefficient map}

Let $X_1,\ldots,X_{q-1}$ be independent variables and define
\begin{equation}\label{eq:formalQ}
 \cQ_q(X)=X_1\sum_{r=1}^{q-1}\beta_rX_r
 -(X_1+X_2)\sum_{r=1}^{q-1}\alpha_rX_r.
\end{equation}
We record its coefficients in the order
\begin{equation}\label{eq:monomialorder}
 X_1^2,\ X_1X_2,\ X_2^2,
 \quad X_1X_r,\ X_2X_r\quad(3\le r\le q-1).
\end{equation}
There are $2q-3$ coordinates.  Write them as
$c_{11},c_{12},c_{22},c_{1r},c_{2r}$.  Direct expansion gives
\begin{align}
 c_{11}&=\beta_1-\alpha_1,\label{eq:c11}\\
 c_{12}&=\beta_2-\alpha_2-\alpha_1,\label{eq:c12}\\
 c_{22}&=-\alpha_2,\label{eq:c22}\\
 c_{1r}&=\beta_r-\alpha_r\qquad(3\le r\le q-1),\label{eq:c1r}\\
 c_{2r}&=-\alpha_r\qquad(3\le r\le q-1).\label{eq:c2r}
\end{align}
Thus we have an integral linear map
\begin{equation}\label{eq:Tq}
 T_q:\Z^{2\pi(q)}\longrightarrow\Z^{2q-3},
 \qquad
 ((a_p)_{p\le q},(b_p)_{p\le q})\longmapsto(c_{ij}).
\end{equation}
The old computations suggested that $T_q$ always has one-dimensional kernel and
that its nonzero Smith invariants are all $1$.  We now prove this uniformly.

\subsection{Uniform kernel, rank, and saturation}

\begin{theorem}[Uniform staircase lattice]\label{thm:uniformstaircase}
Let $q\ge3$ be prime and let $T_q$ be the map in \eqref{eq:Tq}.  Then:
\begin{enumerate}[label=\textup{(\roman*)}]
\item
\[
 \ker T_q
 =\Z\cdot\bigl(e_{a_2}+e_{b_3}\bigr).
\]
In other words, the only formal invisibility is
$a_2\mapsto a_2+t$, $b_3\mapsto b_3+t$, corresponding to
$(M,A)\mapsto(2^tM,3^tA)$.

\item $\rank T_q=2\pi(q)-1$.

\item $\im T_q$ is a saturated sublattice of $\Z^{2q-3}$.

\item Every nonzero invariant factor in the Smith normal form of $T_q$ equals
$1$; explicitly,
\[
 \operatorname{SNF}(T_q)
 =\diag(\underbrace{1,\ldots,1}_{2\pi(q)-1},0,\ldots,0).
\]
\end{enumerate}
\end{theorem}

The proof is constructive.  It gives an integral decoder, which is more useful
than rank alone.

\begin{proof}
We first determine the kernel.  If all coefficients vanish, then
$c_{22}=0$ gives $\alpha_2=0$, hence $a_p=0$ for every $p\ge3$.
For $r\ge3$, equations $c_{2r}=c_{1r}=0$ give $\alpha_r=\beta_r=0$, so
$b_p=0$ for every $p\ge5$.  The remaining equations reduce to
\[
 c_{12}=b_3-a_2=0,
 \qquad
 c_{11}=b_2+b_3-a_2=0.
\]
Thus $b_2=0$ and $b_3=a_2$.  Conversely, adding the same integer $t$ to
$a_2$ and $b_3$ adds $tL_1$ to $\log M$ and $t(L_1+L_2)$ to $\log A$, so
\eqref{eq:formalQ} is unchanged.  This proves (i), and (ii) follows.

For saturation we exhibit quotient coordinates and an integral inverse.  Put
\begin{equation}\label{eq:dcoords}
 d_2=c_{22},\qquad d_r=c_{2r}\ (3\le r\le q-1),\qquad d_q=0.
\end{equation}
For $q\ge5$ also put
\begin{equation}\label{eq:ecoords}
 e_r=c_{1r}-c_{2r}\ (3\le r\le q-1),\qquad e_q=0.
\end{equation}
Equations \eqref{eq:c22} and \eqref{eq:c2r} say $d_r=-\alpha_r$;
therefore, for every prime $p\ge3$,
\begin{equation}\label{eq:adec}
 a_p=d_p-d_{p-1}.
\end{equation}
Similarly $e_r=\beta_r$, and for every prime $p\ge5$,
\begin{equation}\label{eq:bdec}
 b_p=e_{p-1}-e_p.
\end{equation}
The two remaining quotient coordinates are
\begin{equation}\label{eq:b2dec}
 b_2=c_{11}-c_{12}+c_{22}
\end{equation}
and
\begin{equation}\label{eq:deltadec}
 \delta:=b_3-a_2=
 \begin{cases}
 c_{12}-2c_{22},&q=3,\\[0.3em]
 c_{12}-(c_{13}-c_{23})-2c_{22},&q\ge5.
 \end{cases}
\end{equation}
All formulas are integral.  Conversely, from arbitrary integer data
\[
 (a_p)_{3\le p\le q},\quad
 (b_p)_{5\le p\le q},\quad b_2,\quad\delta
\]
choose the gauge $a_2=0$, $b_3=\delta$, reconstruct all tails, and then use
\eqref{eq:c11}--\eqref{eq:c2r}.  This gives an integral section of the quotient
map
\[
 \Z^{2\pi(q)}\longrightarrow
 \Z^{2\pi(q)-1},
 \quad
 (a,b)\longmapsto
 \bigl((a_p)_{p\ge3},(b_p)_{p\ge5},b_2,b_3-a_2\bigr).
\]

Now suppose $k\vct c\in\im T_q$ for a nonzero integer $k$.  Apply the decoder
to $\vct c$.  Since the decoded quotient coordinates of $k\vct c$ are divisible
by $k$, the exponent vector producing $k\vct c$ differs from $k$ times the
integral gauge-fixed reconstruction by a kernel vector.  Applying $T_q$ and
cancelling $k$ in the torsion-free target gives $\vct c\in\im T_q$.  Hence the
image is saturated.  A rank-$r$ integral matrix has saturated image precisely
when all its nonzero Smith invariant factors equal $1$, proving (iii) and (iv).
\end{proof}

\begin{remark}[The $q=3$ endpoint]
When $q=3$ there are four exponent variables and three coefficients.  The same
proof gives rank $3$ and no cokernel torsion.  The special formula in
\eqref{eq:deltadec} is simply the boundary case $\beta_3=0$.
\end{remark}

\subsection{Complete image description}

The decoder yields a local characterization of the image without reference to
prime exponents.

\begin{proposition}[Prime-jump tomography]\label{prop:tomography}
Let $q\ge5$ be prime and let $\vct c\in\Z^{2q-3}$.  Define $d_r,e_r$ by
\eqref{eq:dcoords}--\eqref{eq:ecoords}.  Then $\vct c\in\im T_q$ if and only if
\begin{enumerate}[label=\textup{(\alph*)}]
\item $d_r=d_{r+1}$ whenever $r+1$ is composite and $2\le r\le q-1$;
\item $e_r=e_{r+1}$ whenever $r+1$ is composite and $3\le r\le q-1$;
\item the remaining coordinates $b_2$ and $\delta$ defined by
\eqref{eq:b2dec}--\eqref{eq:deltadec} are integral, which is automatic in the
ambient integer lattice.
\end{enumerate}
For an actual candidate pair, one must add the cone inequalities
\[
 -d_2\ge-d_3\ge\cdots\ge-d_q=0,
 \qquad
 e_3\ge e_4\ge\cdots\ge e_q=0,
\]
and nonnegativity of all reconstructed valuations.
\end{proposition}

\begin{proof}
The tails can jump only at primes, giving the necessity of (a) and (b).  Under
these conditions, \eqref{eq:adec} and \eqref{eq:bdec} define integer prime
valuations, while \eqref{eq:b2dec} and \eqref{eq:deltadec} complete a gauge-fixed
preimage.  The cone conditions are exactly the statement that tail jumps are
nonnegative.
\end{proof}

The term ``tomography'' is literal: a prime exponent is recovered from the
first difference of two adjacent residue coordinates.  No global matrix
inversion is needed.

\subsection{Dimensions and the disappearance of finite-case ambiguity}

For orientation, the relevant dimensions are shown below.
\begin{center}
\begin{tabular}{rrrrr}
\toprule
$q$ & $2\pi(q)$ & $2q-3$ & $\rank T_q$ & $\operatorname{codim}(\im T_q)$\\
\midrule
3  & 4  & 3  & 3  & 0\\
5  & 6  & 7  & 5  & 2\\
7  & 8  & 11 & 7  & 4\\
11 & 10 & 19 & 9  & 10\\
13 & 12 & 23 & 11 & 12\\
17 & 14 & 31 & 13 & 18\\
31 & 22 & 59 & 21 & 38\\
\bottomrule
\end{tabular}
\end{center}
The codimension grows quickly because most adjacent locations are composite and
therefore carry no independent jump.  This sparsity is not numerical accident;
it is the exact prime-boundary geometry of the coefficient lattice.

\subsection{The restricted period relation}

Let $\vct P_q$ denote the vector of real periods in the monomial order
\eqref{eq:monomialorder}:
\begin{equation}\label{eq:Pq}
 \vct P_q=
 \bigl(L_1^2,L_1L_2,L_2^2,(L_1L_r)_{r\ge3},(L_2L_r)_{r\ge3}\bigr).
\end{equation}
Let
\begin{equation}\label{eq:Cqdef}
 \cC_q=\im T_q\subseteq\Z^{2q-3},
 \qquad
 \cC_q^+=T_q(\N_0^{2\pi(q)}).
\end{equation}
Then a $q$-smooth candidate gives
\begin{equation}\label{eq:Cqrelation}
 \vct c\cdot\vct P_q=0,
 \qquad
 \vct c\in\cC_q^+.
\end{equation}
The zero vector corresponds exactly to a control pair $(M,A)=(2^t,3^t)$.
Thus:

\begin{corollary}[Finite smoothness criterion]\label{cor:smoothcriterion}
There is no noncontrol $q$-smooth counterexample if and only if
\[
 \set{\vct c\in\cC_q^+:\vct c\cdot\vct P_q=0}=\set{0}.
\]
Because $\cC_q$ is saturated, any nonzero relation in $\cC_q$ has a primitive
representative in $\cC_q$.
\end{corollary}

The last sentence is operationally important.  A future arithmetic-holonomy
argument may work modulo an auxiliary prime without paying an unnoticed lattice
index: there is no hidden torsion between the exponent lattice modulo controls
and the coefficient lattice.

\section{A fixed-rank unipotent holonomy model}

\subsection{Iterated integrals of the consecutive forms}

For $1\le r\le q-1$ put
\begin{equation}\label{eq:omegar}
 \omega_r=\frac{\dd t}{r+t}
 \qquad(0\le t\le1).
\end{equation}
Then
\begin{equation}\label{eq:intomega}
 \int_0^1\omega_r=L_r.
\end{equation}
Define the two tail forms
\begin{equation}\label{eq:tailforms}
 \Omega_M=\sum_{r=1}^{q-1}\alpha_r\omega_r,
 \qquad
 \Omega_A=\sum_{r=1}^{q-1}\beta_r\omega_r.
\end{equation}
They satisfy
\[
 \int_0^1\Omega_M=\log M,
 \qquad
 \int_0^1\Omega_A=\log A.
\]
For one-forms $\eta,\nu$ on $[0,1]$, write
\begin{equation}\label{eq:iterated}
 I(\eta,\nu)=
 \int_{0<t_1<t_2<1}\eta(t_1)\nu(t_2).
\end{equation}
The shuffle identity gives
\begin{equation}\label{eq:shuffle}
 \left(\int_0^1\eta\right)\left(\int_0^1\nu\right)
 =I(\eta,\nu)+I(\nu,\eta).
\end{equation}
Consequently the determinant \eqref{eq:detzero} equals
\begin{align}
 D(M,A)={}&I(\omega_1,\Omega_A)+I(\Omega_A,\omega_1)\notag\\
 &-I(\omega_1+\omega_2,\Omega_M)
  -I(\Omega_M,\omega_1+\omega_2).
 \label{eq:Diterated}
\end{align}

\subsection{A rank-six connection}

Let $E_{ij}$ be the $6\times6$ matrix unit, with indices $0,1,2,3,4,5$.
Consider the strictly upper triangular matrix-valued one-form
\begin{align}
 \boldsymbol\Omega={}&
 E_{01}\omega_1+E_{15}\Omega_A
 +E_{02}\Omega_A+E_{25}\omega_1\notag\\
 &-E_{03}(\omega_1+\omega_2)+E_{35}\Omega_M
 -E_{04}\Omega_M+E_{45}(\omega_1+\omega_2).
 \label{eq:rank6connection}
\end{align}
The directed graph of nonzero entries has only paths
$0\to\{1,2,3,4\}\to5$.  Hence the path-ordered exponential has no contribution
of length greater than two to its $(0,5)$ entry.

\begin{proposition}[Fixed-rank holonomy realization]\label{prop:rank6}
Let $H$ be the parallel transport of $\dd-\boldsymbol\Omega$ from $0$ to $1$.
Then
\begin{equation}\label{eq:holonomyentry}
 H_{05}=D(M,A)
 =\log2\,\log A-\log3\,\log M.
\end{equation}
The connection is unipotent of rank six and Fuchsian on
\[
 \mathbf P^1\setminus\{-1,-2,\ldots,-(q-1),\infty\}.
\]
Its residue matrices have integer entries.
\end{proposition}

\begin{proof}
The four length-two directed paths contribute, respectively,
\[
 I(\omega_1,\Omega_A),\quad I(\Omega_A,\omega_1),\quad
 -I(\omega_1+\omega_2,\Omega_M),\quad
 -I(\Omega_M,\omega_1+\omega_2).
\]
Equation \eqref{eq:Diterated} proves the first assertion.  Every form
$\omega_r$ has a simple pole at $-r$ with residue $1$, and all coefficients in
\eqref{eq:tailforms} are integers.
\end{proof}

\begin{remark}
The connection rank is independent of $q$.  The difficulty is not a growing
matrix rank but a growing set of rational punctures and unbounded integral
residues.  This is a substantially narrower arithmetic-holonomy target than an
arbitrary weight-two relation among all products $L_iL_j$.
\end{remark}

\subsection{Residue jumps recover prime exponents}

Let $\rho_r^M=\Res_{t=-r}\Omega_M=\alpha_r$ and
$\rho_r^A=\Res_{t=-r}\Omega_A=\beta_r$, with
$\rho_q^M=\rho_q^A=0$.  Then
\begin{equation}\label{eq:resjump}
 a_p=\rho_{p-1}^M-\rho_p^M,
 \qquad
 b_p=\rho_{p-1}^A-\rho_p^A.
\end{equation}
Thus every external valuation is a local jump between two adjacent punctures.
The control direction is the residue gauge transformation
\begin{equation}\label{eq:resgauge}
 (\Omega_M,\Omega_A)\longmapsto
 (\Omega_M+t\omega_1,\ \Omega_A+t(\omega_1+\omega_2)).
\end{equation}
It leaves \eqref{eq:Diterated} unchanged.

This gives a useful restatement of the conjecture.

\begin{target}[Restricted arithmetic holonomy]\label{target:holonomy}
Prove that if the integer residue data of \eqref{eq:rank6connection} have at least
one nonzero prime jump after quotienting by \eqref{eq:resgauge}, then
$H_{05}\ne0$.
\end{target}

Restricted to the nonnegative candidate cone, \cref{target:holonomy} is
equivalent to excluding a $q$-smooth counterexample.  Stated on the full
integer residue lattice it is stronger.  In either form, the formulation matters:
it asks for nonvanishing on a saturated network-flow lattice of residues, not for
linear independence of an arbitrary family of $(2q-3)$ periods.  The strongest currently available
arithmetic-holonomy results prove remarkable near-diagonal cases, but they do
not yet cover this two-star, long-puncture configuration.  The new lattice theorem
shows that there is no further finite-dimensional algebra left to discover at
this interface: any remaining obstruction is genuinely period-theoretic.

\section{The diagonal Kummer supernorm}

\subsection{The elementary determinant cocycle}

Continue to assume a counterexample $x=\beta>1$, and set
\begin{equation}\label{eq:ab}
 a=\log2,\qquad b=\log3.
\end{equation}
For an integer $N\ge2$, choose the positive real roots
\begin{equation}\label{eq:roots}
 P=2^{1/N}=\e^{a/N},\quad
 Q=3^{1/N}=\e^{b/N},\quad
 U=M^{1/N}=\e^{\beta a/N},\quad
 V=A^{1/N}=\e^{\beta b/N}.
\end{equation}
The supplied report studies the algebraic integer
\begin{equation}\label{eq:Xi}
 \Xi_N=(U-1)(Q-1)-(V-1)(P-1).
\end{equation}
Put
\begin{equation}\label{eq:RS}
 R_N=UQ-VP,
 \qquad
 S_N=U+Q-V-P.
\end{equation}
Then
\begin{equation}\label{eq:XiRS}
 \Xi_N=R_N-S_N.
\end{equation}

\begin{proposition}[Third-order distinguished contact]\label{prop:Xiasymp}
As $N\to\infty$,
\begin{equation}\label{eq:Xiasymp}
 \Xi_N=
 \frac{\beta(\beta-1)ab(a-b)}{2N^3}+O_\beta(N^{-4}).
\end{equation}
In particular, $\Xi_N\ne0$ for all sufficiently large $N$.
\end{proposition}

\begin{proof}
Expand every exponential in \eqref{eq:roots}.  The constant and linear terms of
both products in \eqref{eq:Xi} vanish.  The quadratic terms are both
$\beta ab/N^2$ and cancel because $\log M=\beta a$ and
$\log A=\beta b$.  The cubic difference is
\[
 \frac{1}{2N^3}
 \bigl((\beta a)^2b+\beta ab^2-(\beta b)^2a-\beta ba^2\bigr)
 =\frac{\beta(\beta-1)ab(a-b)}{2N^3}.
\]
Since $a\ne b$ and $\beta>1$, the coefficient is nonzero.
\end{proof}

\subsection{Exact common-twist product}

For $\zeta^N=1$, twist all four roots simultaneously and define
\begin{equation}\label{eq:Xizeta}
 \Xi_N(\zeta)=
 (\zeta U-1)(\zeta Q-1)-(\zeta V-1)(\zeta P-1).
\end{equation}
The following identity is the central new Kummer calculation.

\begin{theorem}[Diagonal Kummer norm identity]\label{thm:diagnorm}
For every $N\ge2$,
\begin{equation}\label{eq:diagnorm}
 \Delta_N:=\prod_{\zeta^N=1}\Xi_N(\zeta)
 =R_N^N-S_N^N.
\end{equation}
Consequently $\Delta_N$ is an algebraic integer.  In the full Kummer-rank branch,
\eqref{eq:diagnorm} is the relative norm of $\Xi_N$ along the common diagonal
root-of-unity subgroup.
\end{theorem}

\begin{proof}
Expanding \eqref{eq:Xizeta} gives
\[
 \Xi_N(\zeta)=\zeta^2R_N-\zeta S_N
 =\zeta(\zeta R_N-S_N).
\]
Now
\[
 \prod_{\zeta^N=1}\zeta=(-1)^{N-1}
\]
and
\[
 \prod_{\zeta^N=1}(\zeta R_N-S_N)
 =(-1)^N(S_N^N-R_N^N).
\]
Multiplying the signs yields $R_N^N-S_N^N$.  Algebraic integrality also follows
directly from \eqref{eq:RS}, because $P,Q,U,V$ are algebraic integers.
\end{proof}

A useful degree-reducing form of the same identity is obtained by dividing all
roots by $P$.  Put
\begin{equation}\label{eq:rst}
 r=\frac UP=\left(\frac M2\right)^{1/N},
 \qquad
 s=\frac QP=\left(\frac32\right)^{1/N},
 \qquad
 t=\frac VP=\left(\frac A2\right)^{1/N}.
\end{equation}
Since $P^N=2$,
\begin{equation}\label{eq:reducedDelta}
 \Delta_N=4(rs-t)^N-2(r+s-t-1)^N.
\end{equation}
Thus the common twist has genuinely removed one Kummer direction: $\Delta_N$
lives in a field generated by three radical ratios rather than four independent
roots.

\subsection{Superexponential asymptotic}

The common first two Taylor coefficients of $R_N$ and $S_N$ are
\begin{align}
 c&=(\beta-1)(a-b),\label{eq:cconst}\\
 \gamma&=\frac{(\beta^2-1)(a^2-b^2)}{2}
 =\frac{c(\beta+1)(a+b)}{2}.
 \label{eq:gammaconst}
\end{align}
More precisely,
\begin{equation}\label{eq:RSexpansion}
 R_N=\frac cN+\frac\gamma{N^2}+O_\beta(N^{-3}),
 \qquad
 S_N=\frac cN+\frac\gamma{N^2}+O_\beta(N^{-3}).
\end{equation}
Let
\begin{equation}\label{eq:dconst}
 d=\frac{\beta(\beta-1)ab(a-b)}2.
\end{equation}
Then $R_N-S_N=dN^{-3}+O(N^{-4})$.

\begin{theorem}[Diagonal supernorm asymptotic]\label{thm:supernorm}
For a fixed nonintegral candidate $\beta>1$,
\begin{equation}\label{eq:supernormasymp}
 \abs{\Delta_N}
 \sim
 \abs d\,
 \exp\!\left(\frac{(\beta+1)(a+b)}2\right)
 \abs c^{N-1}N^{-N-1}.
\end{equation}
Equivalently,
\begin{equation}\label{eq:supernormlog}
 \log\abs{\Delta_N}
 =-N\log N+N\log\abs c-\log N+O_\beta(1),
\end{equation}
and
\begin{equation}\label{eq:supernormlimit}
 \lim_{N\to\infty}
 \frac{\log\abs{\Delta_N}+N\log N}{N}
 =\log\bigl((\beta-1)(b-a)\bigr).
\end{equation}
In particular, $\Delta_N$ is nonzero and tends to zero faster than
$C^{-N}$ for every fixed $C>1$ once $N$ is sufficiently large.
\end{theorem}

\begin{proof}
For large $N$, $R_N$ and $S_N$ have the same nonzero sign because
$N R_N,N S_N\to c\ne0$.  The real mean value theorem applied to $z\mapsto z^N$
gives a number $\xi_N$ between $R_N$ and $S_N$ such that
\begin{equation}\label{eq:mvt}
 \Delta_N=N\xi_N^{N-1}(R_N-S_N).
\end{equation}
By \eqref{eq:RSexpansion},
\[
 \xi_N=\frac cN\left(1+\frac{\gamma/c}{N}+O_\beta(N^{-2})\right).
\]
Therefore
\[
 \xi_N^{N-1}
 =\left(\frac cN\right)^{N-1}
 \exp\!\left(\frac\gamma c+o(1)\right).
\]
Using $R_N-S_N=dN^{-3}(1+o(1))$ in \eqref{eq:mvt} gives
\eqref{eq:supernormasymp}; the logarithmic statements follow.
\end{proof}

\begin{remark}[What has actually improved]
The original determinant has only polynomial smallness $N^{-3}$.  The common
orbit product turns the shared first-order scale $c/N$ into $N$ powers while
retaining the third-order difference once.  The gain is therefore the factor
$N^{-N}$, not an unexplained numerical accident.
\end{remark}

\subsection{Numerical calibration and reproducibility}

The asymptotic is numerically stable when $R_N^N-S_N^N$ is evaluated through
\[
 \log|R_N^N-S_N^N|
 =N\log|S_N|+\log\left|\exp\!\left(N\log\frac{|R_N|}{|S_N|}\right)-1\right|,
\]
rather than by direct subtraction of two close floating-point powers.  A
high-precision implementation is included in \cref{app:code}.  Numerical tests
for several real values of $\beta>1$ verify both limits
$N^3\Xi_N\to d$ and \eqref{eq:supernormlimit}.  These tests are calibration only;
the result itself is the exact argument in \cref{thm:supernorm}.

\section{Why the supernorm is not yet a proof}

\subsection{One tiny embedding is inexpensive in a Kummer field}

The element $\Delta_N$ lies in the field generated by the three radical ratios
in \eqref{eq:rst}, so a crude degree bound is $N^3$.  Every conjugate of
$R_N^N-S_N^N$ has absolute value at most $C_\beta^N$ for a constant depending on
the fixed candidate.  The product formula therefore permits one embedding to be
as small as
\begin{equation}\label{eq:liouvillebudget}
 \exp(-O_\beta(N^4)).
\end{equation}
The actual smallness $\exp(-N\log N+O(N))$ is spectacular analytically but is far
above this algebraic lower budget.  A degree collapse by only one Kummer
direction is nowhere near enough.

This can be formulated as a general partial-norm ceiling.

\begin{proposition}[Partial-norm ceiling]\label{prop:partialnormceiling}
Suppose a Kummer family has phase rank $\rho$, and let $H_N$ be a subgroup of
size $N^d$ with $d<\rho$.  Let $\Phi_N$ be an algebraic-integer cocycle whose
complex conjugates are bounded by a fixed constant $C>1$.  If at the
distinguished embedding every factor in the $H_N$-orbit is
$O(N^{-r})$, then the relative norm has size at most
\[
 \exp\bigl(-rN^d\log N+O(N^d)\bigr).
\]
On the other hand, its remaining conjugates can each have size
$C^{N^d}$, and there are on the order of $N^{\rho-d}$ such conjugates.  A bare
product-formula argument therefore allows a negative logarithm of order
$N^\rho$.  Since
\[
 N^d\log N=o(N^\rho)\qquad(d<\rho),
\]
polynomial contact along a proper phase subgroup cannot by itself close the
norm argument.
\end{proposition}

\begin{proof}
The first estimate is the product of $N^d$ factors of size $O(N^{-r})$.  The
house bound for a relative norm is $C^{N^d}$.  Multiplying the distinguished
embedding with the remaining $O(N^{\rho-d})$ embeddings gives the stated
product-formula budget.  The final asymptotic comparison is elementary.
\end{proof}

The proposition isolates what an ideal-divisibility argument would have to do:
it must reduce the generic conjugate budget by an entire power of $N$, not merely
improve the distinguished Taylor order.

\subsection{The root-of-unity phase symbol}

Let $p,q,u,v\in\mathbb S^1$ be independent phases attached to $P,Q,U,V$.
As $N\to\infty$, the corresponding conjugate of \eqref{eq:Xi} tends to
\begin{equation}\label{eq:phasesymbol}
 F(p,q,u,v)=(u-1)(q-1)-(v-1)(p-1).
\end{equation}
Expanded as a Laurent polynomial,
\begin{equation}\label{eq:Fexpanded}
 F=uq-vp-u-q+v+p.
\end{equation}
For the diagonal product, the two leading coefficients are
\begin{equation}\label{eq:r0s0}
 r_0=uq-vp,
 \qquad
 s_0=u+q-v-p.
\end{equation}
The simultaneous zero locus of these coefficients is completely elementary but
important.

\begin{lemma}[Two-axis simultaneous-zero locus]\label{lem:twoaxis}
For nonzero complex numbers $p,q,u,v$, the equations
\[
 uq=vp,
 \qquad
 u+q=v+p
\]
hold if and only if one of the following alternatives holds:
\begin{equation}\label{eq:axes}
 u=p,\ v=q,
 \qquad\text{or}\qquad
 q=p,\ v=u.
\end{equation}
\end{lemma}

\begin{proof}
Substitute $v=u+q-p$ into $uq=vp$.  One obtains
\[
 0=uq-p(u+q-p)=(u-p)(q-p).
\]
If $u=p$, then $vp=uq=pq$ gives $v=q$.  If $q=p$, then $vp=up$ gives $v=u$.
The converse is immediate.
\end{proof}

Thus only two natural phase axes make both bases of the diagonal difference
start at order $N^{-1}$.  The distinguished common diagonal is their
intersection.  This explains geometrically why a second relative norm does not
automatically repeat the first gain: most quotient phases leave at least one of
$r_0,s_0$ nonzero.

\subsection{Mahler measure and positive phase entropy}

For a Laurent polynomial $G$ in $m$ variables, write
\begin{equation}\label{eq:mahlerdef}
 m(G)=\frac1{(2\pi)^m}
 \int_{[0,2\pi]^m}\log\abs{G(\e^{i\theta_1},\ldots,\e^{i\theta_m})}
 \dd\theta_1\cdots\dd\theta_m
\end{equation}
for its logarithmic Mahler measure.  The multivariable Kronecker theorem~\cite{Boyd1981,Lawton1983,Smyth1981} says
that a primitive integral Laurent polynomial has $m(G)\ge0$, with equality only
for a monomial times a product of cyclotomic polynomials in monomials.

\begin{theorem}[Positive phase entropy]\label{thm:positiveMahler}
The phase polynomial $F$ in \eqref{eq:phasesymbol} has
\begin{equation}\label{eq:mFpositive}
 m(F)>0.
\end{equation}
Numerically,
\begin{equation}\label{eq:mFnumeric}
 m(F)\approx0.54807.
\end{equation}
\end{theorem}

\begin{proof}
Regard $F$ as a polynomial in $u$ over $\Q[p,q,v]$:
\[
 F=(q-1)u-\bigl(q-1+(v-1)(p-1)\bigr).
\]
The two coefficients are coprime, so Gauss's lemma shows that $F$ is irreducible.
Its Newton polytope is not a line segment: the six exponent vectors in
\eqref{eq:Fexpanded} have affine span of dimension greater than one.  If
$m(F)=0$, multivariable Kronecker would force the irreducible polynomial $F$ to
be cyclotomic in a single monomial, whose Newton polytope is a line segment.
This contradiction proves strict positivity.

For the numerical value, integrate first in $v$.  Jensen's formula gives
\begin{equation}\label{eq:mFjensen}
 m(F)=\frac1{(2\pi)^3}\int
 \log\max\bigl(\abs{uq-u-q+p},\abs{1-p}\bigr)
 \dd\theta_u\dd\theta_q\dd\theta_p.
\end{equation}
A scrambled Sobol computation with $2^{22}$ points gives the value in
\eqref{eq:mFnumeric}; code appears in \cref{app:code}.
\end{proof}

\begin{remark}[What strict positivity means]
The geometric mean of the leading conjugate magnitudes is
$\exp(m(F))>1$.  Therefore a full root-of-unity norm of the unnormalized simple
cocycle is phase-expansive at leading order.  The sparse set of very small
conjugates cannot be treated as representative of the full orbit.
\end{remark}

\subsection{Cyclotomic units cannot erase the entropy}

Kummer fields contain many units whose distinguished absolute value is tiny;
for example, differences of radicals can have norm $1$ in a suitable relative
extension.  Multiplying by such a unit can make one embedding arbitrarily
smaller while making other embeddings larger.  At the phase-symbol level these
normalizations multiply $F$ by generalized cyclotomic Laurent polynomials, all
of Mahler measure zero.  Hence
\begin{equation}\label{eq:unitinvariance}
 m(FG)=m(F)+m(G)=m(F)>0
\end{equation}
for every generalized cyclotomic $G$.

This is why the Mahler calculation is stronger than a crude conjugate bound.
It is invariant under the obvious unit rebalancings.  A successor cocycle must
alter the noncyclotomic phase symbol itself.

\begin{criterion}[Zero-entropy Kummer design criterion]\label{crit:zeroentropy}
A plausible full-norm Kummer proof should construct algebraic integers
$\Phi_N$ satisfying all of the following.
\begin{enumerate}[label=\textup{(\roman*)}]
\item At the distinguished embedding, $\Phi_N$ has high-order contact forced by
\eqref{eq:detzero}.
\item Its leading root-of-unity phase symbol $G$ has $m(G)=0$, hence is
generalized cyclotomic.
\item After removing the cyclotomic leading term, the first nonzero descendant is
control-sensitive: in staircase coordinates it contains a nonzero prime jump
$a_p$ or $b_p$ with $p\ge5$.
\item The required normalization is integral or is backed by a certified ideal
divisibility whose norm pays for every denominator.
\end{enumerate}
\end{criterion}

Condition (ii) suppresses the phase entropy; condition (iii) prevents the
construction from collapsing to a tautology on the control pair.  The supplied
report had versions of (i) and (iv) separately.  The new Mahler audit shows that
(ii) and (iii) must be enforced simultaneously.

\section{A unifying two-star viewpoint}

\subsection{The same determinant appears twice}

The staircase period polynomial is
\begin{equation}\label{eq:twostarperiod}
 X_1B(X)-(X_1+X_2)A(X),
\end{equation}
where $A$ and $B$ are tail-linear forms.  The Kummer phase symbol is
\begin{equation}\label{eq:twostarphase}
 (u-1)(q-1)-(v-1)(p-1).
\end{equation}
Both are rank-two determinants with two distinguished structural directions and
two output directions.  In the period model, taking discrete residue differences
recovers external prime exponents.  In the Kummer model, taking a common phase
norm amplifies the distinguished cancellation but exposes a positive-entropy
torus polynomial.

This suggests that the next construction should not be sought as an unrelated
large determinant.  It should be a \emph{filtered two-star cocycle}: the leading
phase layer should factor cyclotomically, while the next layer should equal a
prime-jump functional from \cref{prop:tomography}.

\subsection{A concrete hybrid ansatz}

Let $\vct c\in\cC_q$ be a primitive staircase coefficient vector and let
$J_p(\vct c)$ denote one of the decoded jumps in
\eqref{eq:adec}--\eqref{eq:bdec}.  A useful hybrid construction would be a family
of Laurent polynomials
\begin{equation}\label{eq:hybridPhi}
 \Phi_{N,\vct c}(\vct z;N^{-1})
 =G_0(\vct z)^{N}
 +\frac1N G_{1,\vct c}(\vct z)
 +\frac1{N^2}G_{2,\vct c}(\vct z)+\cdots
\end{equation}
with the following properties:
\begin{itemize}
\item $G_0$ is generalized cyclotomic;
\item the first coefficient not divisible by $G_0$ contains
$J_p(\vct c)$ for some external prime $p$;
\item evaluation at the positive radical point is an algebraic integer and has
superexponential decay;
\item the degree and coefficient height grow slowly enough for a full norm.
\end{itemize}
The ansatz is intentionally algebraic.  It translates the transcendence problem
into a finite symbolic search at each support size.  The uniform decoder makes
control sensitivity machine-checkable: a candidate cocycle that annihilates all
$J_p$ can be rejected before any analytic estimate.

\subsection{Why a larger Taylor order alone is insufficient}

Suppose a cocycle has order-$r$ contact along a $d$-dimensional phase subgroup.
Its partial norm gains $rN^d\log N$ in the negative logarithm.  If $d$ is smaller
than the ambient phase rank, \cref{prop:partialnormceiling} shows that increasing
$r$ by any fixed amount does not change the exponent of $N$ and therefore cannot
beat the remaining conjugates.  A successful construction must do at least one
of the following:
\begin{enumerate}[label=\textup{(\alph*)}]
\item make the low-order cancellation valid on the full relevant phase torus;
\item force a degree collapse to $N^{d+o(1)}$;
\item divide by a genuinely large common ideal whose norm cancels the generic
conjugate budget;
\item or replace a norm argument by a trace or matrix coefficient with certified
sign/cancellation control.
\end{enumerate}
The Mahler criterion is the cleanest version of (a).

\section{Further directions audited in this continuation}

\subsection{Restricted Calegari--Dimitrov--Tang input}

The recent arithmetic-holonomy results of Calegari--Dimitrov--Tang~\cite{CDT2024,CDTsurvey2026} prove
linear independence for striking families of logarithmic periods, including
near-diagonal consecutive logarithms.  The staircase relation here is different:
it uses the fixed hubs $L_1,L_2$ and all tails $L_r$.  The coefficient lattice is
far smaller than the full integer lattice, but its coefficients are unbounded.

The uniform theorem suggests a realistic specialization of the general theory.
One does not need arbitrary linear independence of all products in
\eqref{eq:Pq}.  It is enough to prove nonvanishing on the network lattice
$\cC_q$, or even on its nonnegative cone $\cC_q^+$.  A proof may exploit:
\begin{itemize}
\item the consecutive poles $-1,-2,\ldots,-(q-1)$;
\item the tail monotonicity of actual candidates;
\item the fact that every noncontrol vector has a nonzero prime jump;
\item the rank-six realization \eqref{eq:rank6connection};
\item and the absence of cokernel torsion.
\end{itemize}
No theorem presently available to the author establishes this restricted
nonvanishing, but the target is now exact enough to formalize and test.

\subsection{Rank-three multiplicative branch}

Let $\rho$ be the multiplicative rank of $2,3,M,A$.  The supplied report reduces
the difficult Kummer geometry to $\rho\in\{3,4\}$.  In rank four, all independent
root-of-unity phases are present, so \cref{thm:positiveMahler} directly describes
the leading full phase space.  In rank three, the phases lie on a codimension-one
subtorus determined by a primitive multiplicative relation
\begin{equation}\label{eq:rank3relation}
 2^r3^sM^uA^v=1.
\end{equation}
The leading symbol becomes the restriction of $F$ to that subtorus.

A potentially decisive classification problem is therefore:

\begin{problem}[Rank-three zero-entropy classification]\label{prob:rank3entropy}
Classify primitive $(r,s,u,v)\in\Z^4$ for which the restriction of
$F=(u_0-1)(q_0-1)-(v_0-1)(p_0-1)$ to the subtorus
$p_0^rq_0^su_0^uv_0^v=1$ has Mahler measure zero.  Show, if true, that every such
case is incompatible with a noncontrol candidate.
\end{problem}

This is a finite Newton-polytope and cyclotomic-factorization problem before any
transcendence estimate is used.  It appears more tractable than a direct analysis
of all rank-three Kummer conjugates.  I have not completed the classification,
so it is recorded as a target rather than a theorem.

\subsection{The Sturmian/automatic route}

The carry word
\begin{equation}\label{eq:carry}
 c_n=\floor{(n+1)x}-\floor{nx}
\end{equation}
is Sturmian for nonintegral $x$.  The integer recurrences
\[
 2^{\{(n+1)x\}}=\frac{M}{2^{c_n}}2^{\{nx\}},
 \qquad
 3^{\{(n+1)x\}}=\frac{A}{3^{c_n}}3^{\{nx\}}
\]
realize the same word in dyadic and triadic arithmetic.  A proof that this word
were both $2$-automatic and $3$-automatic would finish the problem by Cobham's
theorem, because an aperiodic Sturmian word is not automatic.  The missing step
is still missing: integrality of the two recurrences does not by itself provide
a finite automaton reading the index $n$.  The state variables are rational but
their denominator exponents grow without bound.  No honest finite-state
compression was found in this continuation.

\subsection{Small-type entire interpolation}

The supplied report identifies a powerful sufficient condition: a nonconstant
entire function of sufficiently small positive quadratic type, integer-valued on
$\N_0+x\N_0$, would contradict a counterexample through a Pila-style counting
argument.  Canonical products that vanish on the semigroup have infinite
quadratic type, and standard Barnes double-gamma or double-sine constructions
inherit the same boundary-density cost.  I examined whether one could avoid
zeros and impose only integer values through a rapidly convergent Newton series.
The obstruction remains the inverse of the near-collision Vandermonde matrix:
small gaps in $\N_0+x\N_0$ force coefficient growth at the same critical order.
No subcritical construction was obtained.

The new Kummer supernorm does, however, suggest an interpolation analogue:
replace exact cardinal functions by orbit products whose leading potential has
zero logarithmic energy.  In the entire-function language this is again a
zero-Mahler/zero-capacity condition.  The parallel is conceptually useful, but it
has not yet produced a function meeting the required type bound.

\subsection{Prime distribution and external descendants}

The decoder \eqref{eq:adec}--\eqref{eq:bdec} means that the largest external
prime need not be tracked through a global factorization.  It is a boundary jump
of the tail connection.  A possible use of prime distribution is therefore not
to locate a prime in a short interval, but to produce many long composite
stretches on which residues are constant and then compare holonomy across their
endpoints.  Existing prime-gap bounds are too weak by themselves, because the
residue heights can be arbitrarily large.  A useful theorem would need an
output-sensitive estimate whose gain depends on the nonzero jump and whose loss
depends only logarithmically on the tail height.

\section{Sharply stated next theorem targets}

\subsection{Target A: two-star holonomy nonvanishing}

\begin{target}[Saturated two-star nonvanishing]\label{target:A}
For every prime $q\ge5$ and every nonzero primitive
$\vct c\in\cC_q$, prove
\[
 \vct c\cdot\vct P_q\ne0.
\]
A weaker but still sufficient version may restrict to the cone $\cC_q^+$.
\end{target}

This is the most direct continuation of current arithmetic-holonomy technology.
The lattice theorem removes all algebraic ambiguity from its statement.

\subsection{Target B: cyclotomic leading Kummer symbol}

\begin{target}[Control-sensitive zero-entropy cocycle]\label{target:B}
Construct an algebraic-integer Kummer cocycle whose leading phase symbol is a
product of cyclotomic polynomials in monomials, while its first noncyclotomic
Taylor coefficient is a nonzero decoded prime jump from
\cref{prop:tomography}.
\end{target}

The requirement that the first surviving coefficient be a decoded jump is
stronger than merely asking the cocycle to be nonzero at a generic point.  It
prevents a construction from detecting only the control shift.

\subsection{Target C: rank-three entropy classification}

\begin{target}[Rank-three restricted Mahler positivity]\label{target:C}
Prove that for every rank-three relation compatible with a noncontrol candidate,
the restriction of $F$ in \eqref{eq:phasesymbol} to the corresponding Kummer
subtorus has positive Mahler measure.
\end{target}

This would close the leading-symbol audit in both possible multiplicative-rank
branches.  It would not prove \AEconj by itself, but it would rule out an entire class
of low-order Kummer norms uniformly.

\subsection{Target D: ideal normalization with the correct exponent}

\begin{target}[Norm-paying common ideal]\label{target:D}
For a zero-entropy cocycle as in \cref{target:B}, find a common ideal
$\mathfrak d_N$ dividing all phase conjugates such that
\[
 \log\Norm(\mathfrak d_N)
\]
pays the full generic conjugate budget, of order $N^\rho$, while division by
$\mathfrak d_N$ preserves a nonzero external-prime descendant.
\end{target}

The exponent is essential.  Divisibility of size $\exp(O(N^{\rho-1}\log N))$
will not suffice against $\exp(O(N^\rho))$ generic conjugates.

\subsection{Priority ordering}

The recommended order of attack is:
\begin{enumerate}
\item Formalize \cref{thm:uniformstaircase,prop:rank6,thm:diagnorm,lem:twoaxis}.
This gives a trusted algebraic substrate and prevents future computational
experiments from drifting into the control kernel.
\item Search symbolically for low-support cocycles satisfying
\cref{target:B}.  Reject candidates by Mahler measure and the prime-jump decoder
before performing any field-norm calculation.
\item In parallel, ask whether the arithmetic-holonomy machinery can exploit the
rank-six connection and monotone tail cone, rather than arbitrary coefficient
vectors.
\item Treat the rank-three Mahler classification as a finite combinatorial
subproject.  Its output will clarify whether rank three offers a genuinely softer
Kummer branch.
\end{enumerate}

\section{Lean formalization blueprint}

The new exact statements were chosen to separate cleanly into finite algebra and
analysis.  The following module plan should be practical in Mathlib.

\subsection{Module \texttt{StaircaseTail}}

Define finite prime-indexed exponent functions and tails
\[
 \alpha(r)=\sum_{p\in\cP_q,\ r<p}a(p).
\]
Prove the finite-sum jump lemma
\[
 \alpha(p-1)-\alpha(p)=a(p)
\]
for prime $p$, and constancy across composite boundaries.  This module is purely
combinatorial.

\subsection{Module \texttt{StaircaseCoefficient}}

Define $T_q$ by the explicit formulas
\eqref{eq:c11}--\eqref{eq:c2r}.  Formalize:
\begin{lstlisting}
theorem kernel_Tq : LinearMap.ker (Tq q) =
  Submodule.span Z {controlVector q}

 theorem rank_Tq : finrank Z (LinearMap.range (Tq q)) = 2 * pi q - 1
\end{lstlisting}
Rather than formalizing Smith normal form first, prove saturation from the
explicit decoder and section.  Smith invariants can be added as a corollary.

\subsection{Module \texttt{RankSixPathAlgebra}}

Represent the directed acyclic graph with four length-two paths.  The identity
$H_{05}=D(M,A)$ follows from finite path enumeration and the shuffle identity.
One need not formalize analytic differential equations initially; a formal Chen
series in a noncommutative path algebra is enough.

\subsection{Module \texttt{DiagonalNormIdentity}}

The key polynomial identity is independent of transcendence:
\begin{lstlisting}
theorem prod_roots_mul_sub (R S : C) (N : N) :
  (prod z in rootsOfUnity N, z * (z * R - S)) = R^N - S^N
\end{lstlisting}
In practice it is cleaner to prove the universal polynomial identity using
$X^N-Y^N$ and then specialize to roots of unity.

\subsection{Module \texttt{PhaseAxes}}

The two-axis lemma is a short integral-domain proof:
\begin{lstlisting}
theorem phase_axes
  (hp : p != 0) (h1 : u*q = v*p) (h2 : u+q = v+p) :
  (u = p /\ v = q) \/ (q = p /\ v = u)
\end{lstlisting}
The factorization $(u-p)(q-p)=0$ is the core step.

\subsection{Analytic asymptotics}

The asymptotic theorem can be postponed until the algebraic identities are
verified.  A first formal milestone is to prove the limits
\[
 NR_N\to c,\quad N^2(R_N-c/N)\to\gamma,
 \quad N^3(R_N-S_N)\to d.
\]
The mean-value proof of \cref{thm:supernorm} then uses real analysis already
available in Mathlib, although managing $N$-dependent powers will require some
auxiliary logarithm lemmas.

\section{Conclusions}

The work in this continuation changes the shape of the problem in two concrete
ways.

On the period side, there is no longer an unresolved finite-rank pattern.  The
consecutive-log coefficient map is completely understood for every cutoff.  It
is a split, saturated encoding of all prime valuations modulo exactly one
control line.  Every external prime is a local residue jump, and the complete
period obstruction is one matrix coefficient of a rank-six unipotent
connection.

On the Kummer side, the elementary determinant is more powerful than its
$N^{-3}$ Taylor estimate suggests.  Its common diagonal orbit is an explicit
nonzero algebraic integer of size $\exp(-N\log N+O(N))$.  The failure of the norm
route is correspondingly more informative: it is not caused by insufficient
smallness at the distinguished place, but by positive Mahler entropy of the
leading phase symbol.  Cyclotomic units cannot remove that entropy.

The strongest honest statement at present is therefore not a proof but a
sharp bifurcation.  A successful proof must either establish restricted
arithmetic-holonomy nonvanishing on the saturated two-star lattice, or construct
a new Kummer cocycle with cyclotomic leading phase and an external-prime first
descendant.  Both targets are finite enough to support symbolic search and Lean
formalization, and both expose exactly where any purported breakthrough must do
something unavailable in the current report.

\appendix

\section{Explicit small-cutoff matrices}\label{app:matrices}

For $q=3$, order the exponent variables as
$(a_2,a_3,b_2,b_3)$ and coefficients as $(c_{11},c_{12},c_{22})$.  Then
\[
 T_3=
 \begin{pmatrix}
 -1&-1&1&1\\
 -1&-2&0&1\\
 0&-1&0&0
 \end{pmatrix}.
\]
Its kernel is generated by $(1,0,0,1)$, and its determinant minors have gcd $1$.

For $q=5$, order variables as
$(a_2,a_3,a_5,b_2,b_3,b_5)$ and coefficients as
$(c_{11},c_{12},c_{22},c_{13},c_{23},c_{14},c_{24})$.  The tail formulas give
\[
 T_5=
 \begin{pmatrix}
 -1&-1&-1&1&1&1\\
 -1&-2&-2&0&1&1\\
 0&-1&-1&0&0&0\\
 0&0&-1&0&0&1\\
 0&0&-1&0&0&0\\
 0&0&-1&0&0&1\\
 0&0&-1&0&0&0
 \end{pmatrix}.
\]
The repeated rows reflect the composite boundary $4$: tails cannot jump between
$r=3$ and $r=4$.  The rank is $5$, the kernel is again the control line, and the
image has codimension $2$ with no torsion.

\section{Taylor details for the Kummer supernorm}\label{app:taylor}

Set $t=1/N$.  Then
\begin{align*}
 R_N&=\e^{(\beta a+b)t}-\e^{(\beta b+a)t},\\
 S_N&=\e^{\beta at}+\e^{bt}-\e^{\beta bt}-\e^{at}.
\end{align*}
The first coefficient of each is
\[
 (\beta a+b)-(\beta b+a)=(\beta-1)(a-b)=c.
\]
The second coefficient of $R_N$ is
\[
 \frac{(\beta a+b)^2-(\beta b+a)^2}{2}
 =\frac{(\beta^2-1)(a^2-b^2)}2=\gamma,
\]
and the second coefficient of $S_N$ is the same.  Their cubic difference is
\begin{align*}
 6[t^3](R_N-S_N)
 ={}&(\beta a+b)^3-(\beta b+a)^3\\
 &-(\beta a)^3-b^3+(\beta b)^3+a^3\\
 ={}&3\beta(\beta-1)ab(a-b),
\end{align*}
which gives $d$ in \eqref{eq:dconst}.  The refined asymptotic constant in
\eqref{eq:supernormasymp} comes from
\[
 \left(1+\frac{\gamma/c}{N}+O(N^{-2})\right)^{N-1}
 \longrightarrow\e^{\gamma/c}.
\]

\section{Reproducible computations}\label{app:code}

The following Python code estimates $m(F)$ after the exact Jensen reduction
\eqref{eq:mFjensen}.  It requires NumPy and SciPy.

\begin{lstlisting}[language=Python]
import numpy as np
from scipy.stats import qmc


def mahler_F(power=22, seed=12345):
    sampler = qmc.Sobol(d=3, scramble=True, seed=seed)
    x = sampler.random_base2(m=power)
    theta = 2.0 * np.pi * x

    u = np.exp(1j * theta[:, 0])
    q = np.exp(1j * theta[:, 1])
    p = np.exp(1j * theta[:, 2])

    a = u * q - u - q + p
    b = 1.0 - p
    values = np.log(np.maximum(np.abs(a), np.abs(b)))
    return float(values.mean())


print(mahler_F())  # approximately 0.54807
\end{lstlisting}

The next routine checks the diagonal asymptotic with arbitrary real $\beta>1$.
It is an analytic calibration and does not assume that the chosen $\beta$ gives
integer powers.

\begin{lstlisting}[language=Python]
import mpmath as mp


def diagonal_data(beta, N, digits=200):
    mp.mp.dps = digits
    beta = mp.mpf(beta)
    N = int(N)
    a, b = mp.log(2), mp.log(3)

    P = mp.e ** (a / N)
    Q = mp.e ** (b / N)
    U = mp.e ** (beta * a / N)
    V = mp.e ** (beta * b / N)

    R = U * Q - V * P
    S = U + Q - V - P
    Xi = R - S

    # Stable evaluation for large N, when R and S have the same sign.
    z = N * mp.log(abs(R) / abs(S))
    log_abs_Delta = N * mp.log(abs(S)) + mp.log(abs(mp.expm1(z)))

    c = (beta - 1) * (a - b)
    renormalized = (log_abs_Delta + N * mp.log(N)) / N
    return N**3 * Xi, log_abs_Delta, renormalized, mp.log(abs(c))
\end{lstlisting}

For the staircase theorem, direct matrix construction is useful only as a test;
the proof is the decoder in \cref{thm:uniformstaircase}.  A short SymPy check is:

\begin{lstlisting}[language=Python]
import sympy as sp
from sympy import primerange
from sympy.matrices.normalforms import smith_normal_form
from sympy.polys.domains import ZZ


def staircase_matrix(q):
    primes = list(primerange(2, q + 1))
    columns = []

    for side in (0, 1):       # 0 = M exponents, 1 = A exponents
        for p in primes:
            a = {r: int(side == 0 and p > r) for r in range(1, q)}
            b = {r: int(side == 1 and p > r) for r in range(1, q)}
            col = [b[1] - a[1], b[2] - a[2] - a[1], -a[2]]
            for r in range(3, q):
                col += [b[r] - a[r], -a[r]]
            columns.append(col)

    return sp.Matrix.hstack(*[sp.Matrix(c) for c in columns])


for q in (3, 5, 7, 11, 13):
    T = staircase_matrix(q)
    S = smith_normal_form(T, domain=ZZ)
    nonzero = [S[i, i] for i in range(min(S.shape)) if S[i, i] != 0]
    print(q, T.rank(), nonzero)
\end{lstlisting}

\section{Selected references}

\begin{thebibliography}{99}

\bibitem{AE1944}
L.~Alaoglu and P.~Erd\H{o}s,
\newblock On highly composite and similar numbers,
\newblock \emph{Transactions of the American Mathematical Society} 56 (1944),
448--469.

\bibitem{Baker1975}
A.~Baker,
\newblock \emph{Transcendental Number Theory},
\newblock Cambridge University Press, 1975.

\bibitem{Boyd1981}
D.~W.~Boyd,
\newblock Kronecker's theorem and Lehmer's problem for polynomials in several
variables,
\newblock \emph{Journal of Number Theory} 13 (1981), 116--121.

\bibitem{CDT2024}
F.~Calegari, V.~Dimitrov, and Y.~Tang,
\newblock The linear independence of $1$, $\zeta(2)$, and
$L(2,\chi_{-3})$,
\newblock arXiv:2408.15403, 2024--2026.

\bibitem{CDTsurvey2026}
F.~Calegari, V.~Dimitrov, and Y.~Tang,
\newblock Arithmetic holonomy bounds and effective Diophantine approximation,
\newblock ICM survey, arXiv:2510.04156, revised 2026.

\bibitem{DasguptaKakdeII}
S.~Dasgupta and M.~Kakde,
\newblock On the ranks of matrices of logarithms of algebraic numbers, II,
\newblock \emph{Advances in Mathematics}, 2026; arXiv:2408.08178.

\bibitem{Cobham1969}
A.~Cobham,
\newblock On the base-dependence of sets of numbers recognizable by finite automata,
\newblock \emph{Mathematical Systems Theory} 3 (1969), 186--192.

\bibitem{Lawton1983}
W.~M.~Lawton,
\newblock A problem of Boyd concerning geometric means of polynomials,
\newblock \emph{Journal of Number Theory} 16 (1983), 356--362.

\bibitem{Smyth1981}
C.~J.~Smyth,
\newblock On measures of polynomials in several variables,
\newblock \emph{Bulletin of the Australian Mathematical Society} 23 (1981),
49--63.

\bibitem{Unified6}
\emph{Alaoglu--Erd\H{o}s Unified Report 6},
\newblock supplied research report, 22--24 August 2026.

\bibitem{Waldschmidt}
M.~Waldschmidt,
\newblock The four exponentials problem and Schanuel's conjecture,
\newblock in \emph{Mathematics Going Forward}, Lecture Notes in Mathematics,
Springer, 2023, 579--592.

\end{thebibliography}

\end{document}
